% s-195-p195-m002; original PDF 195, 196
\Aphorism{s-195-p195-m002}{282}
Wir treiben \Sperr{Intimolatrie} — Verzeihung für diesen graecolatinischen Wortcentauren, der es mit „Sociologie“ und „Phänomenalismus“ aufnehmen kann. Wir haben das Binnenleben der Seele entdeckt und freuen uns als echte Künstler, Deutsche und Obscuranten, dass es wieder einmal eine Qualität giebt, die sich nicht beweisen, sondern nur glauben lässt. Intim! ein neues Evangelium, eine neue Rangordnung und Gnadenwahl: und die Ersten werden die Letzten sein \Dots{3} Der Sprachbegabte, der Formvollendete, der Werkschöpfer — das sind arme Oberflächler; wer aber schweigt, oder stammelt, oder delirirt, der ist intim, der führt ein Binnenleben. Worte finden ist gemein; das Unaussprechliche \Dots{3} die versagende Stimme \Dots{3} das ist \Dots{3} das Höchste \Dots{3} Gedanken und Anschauung sind antiquirt; wir wollen „Stimmung“, die Dinge „betrachtet in ihrer Unlösbarkeit vom Milieu“. Denn: ein Gedanke lässt sich wägen, abgrenzen, mittheilen, aber Stimmung ist das Undefinirbare: wir schlagen sie an, in dir klingt sie mit, oder du bist nicht intim, du hast kein Binnenleben. Goethe ist nicht intim, denn er konnte reden, bilden; das Ungeheure schnürte ihm nicht die Kehle zu, wie dem titanischen Nichtskönner Grabbe. Hebbel, ein Gewaltiger; aber seine Dramen und Gedichte sind nichts; die Tagebücher allein sind etwas, denn — sie sind ungeformt, was die Dramen und Gedichte leider nicht sind. Das Bewusstsein — langweilig, philisterhaft; Traum, Hypnose, Ekstase sind die eigentlichen Emanationen des Binnenlebens; Helmholtz weiss nichts, Sar Peladan weiss alles. Bethätigung, Wissenschaft, Erfolg — alles keine Intimität; nur der schlafende Fakir taucht in die Tiefen, des Idioten weisse Seele ist der Sitz eleusinischer Enthüllungen. —— Nun, auch diese Pendelschwingung der modernen Entwicklung ist schon wieder in der Umkehr begriffen. Von den Pferdekräften zur Seelenmagie, von der Mechanik zur Mystik, vom Militarismus zum Occultismus — und wieder zurück; aus dem Hin und Her wird doch schliesslich einmal ein Vorwärts.
\par

